
\section{Introduction}
\label{sec:introduction}
Vehicles are increasingly becoming software cyber-physical systems in which perception, control, communication, and user-facing services must operate under strict resource, timing, thermal, and dependability constraints. Artificial intelligence is now used in advanced driver-assistance systems (ADAS) for tasks such as traffic sign recognition, lane-line detection, driver monitoring, and sensor fusion. Although neural-network classifiers can achieve high predictive accuracy, predictions correct on a finite test set may still change under small variations in image or sensor values. Formal local-robustness analysis can therefore provide guarantees over explicitly defined input regions, complementing conventional accuracy evaluation.

Deploying neural networks on embedded automotive platforms also requires reducing their processing and memory needs. Quantization is a common approach because integer-only inference can reduce model size and memory traffic while exploiting efficient integer arithmetic~\cite{Jacob2018Integer}. However, quantization alters the network's numerical semantics. Products must be rescaled, sums are accumulated with finite width, values are rounded, and intermediate activations may be saturated or clamped. Consequently, a robustness property established for a real-valued network does not automatically hold for the fixed-point implementation executed on the target platform.

This distinction becomes more relevant when the neural network is integrated into a complete software environment such as \AAOS{}. The deployed artifact includes not only the quantized network and its input encoder, but also the generated C implementation, the native executable, the compiler, the runtime, and the hardware configuration. A deployment-aware analysis should therefore distinguish guarantees established for an abstract neural-network model from guarantees established for the finite-width program that performs inference.

Present methods address different parts of this problem. Abstract-domain methods such as DeepPoly provide scalable bounds for neural-network verification~\cite{Singh2019DeepPoly}. QNNVerifier translates quantized neural-network implementations into SMT-based software-verification problems and models finite word-length arithmetic~\cite{Song2021QNNVerifier}. CEG4N combines search-based quantization with equivalence checking~\cite{Matos2022CEG4N}, while Quadapter synthesizes layer-wise quantization strategies using backward preimages and MILP under-approximations~\cite{Zhang2024Quadapter}. These approaches demonstrate the value of formal analysis regarding quantized neural networks, but the relationship among quantization, generated fixed-point C code, formal verification, and execution on an automotive Android platform remains insufficiently characterized.

In this work, we study this relationship using \PreQBMC{}. The workflow connects fixed-point quantization, generated C arithmetic, local-robustness contracts, and deployment on an \AAOS{} platform. Rather than treating a vendor-specific accelerator or NPU kernel as formally verified, we use the generated fixed-point C implementation as the verification boundary. \ESBMC{} checks generated certificate obligations linked to the C implementation, while the same native implementation is executed on a Khadas VIM3. This design separates the formal verification artifact from the native execution experiment. The main contributions of this work are:

\begin{enumerate}[leftmargin=*,nosep]
  \item a deployment-aware methodology that connects fixed-point quantization, explicit C arithmetic, and local-robustness contracts;
  \item a certification criterion covering input encoding, layer obligations, contract composition, output margin, and stated arithmetic-safety assumptions; and
  \item native execution on an \AAOS{}-based board and a single-threaded \Aseven{} evaluation using the same reference C source.
\end{enumerate}

The remainder of this paper is organized as follows. Section~\ref{sec:related} reviews work on quantization and neural-network verification. Section~\ref{sec:methodology} describes the deployment-aware methodology, including fixed-point semantics, verification contracts, and \AAOS{} integration. Section~\ref{sec:setup} presents the research questions and experimental arrangement. Section~\ref{sec:results} reports the verification and deployment results, while Section~\ref{sec:limitations} discusses limitations and threats to validity. Finally, Section~\ref{sec:conclusion} concludes the paper with the contributions.


\section{Related Work}
\label{sec:related}

\subsection{Fixed-Point CNNs for Embedded Inference}



Fixed-point quantization can reduce the memory and computational requirements of CNNs on embedded platforms. Goyal \emph{et al.} select layer-wise fixed-point formats while controlling accuracy loss~\cite{Goyal2021FixedPointCNN}, whereas Romano \emph{et al.} demonstrate the importance of evaluating integer inference directly on resource-constrained hardware~\cite{Romano2026NeuroC} [AUTHOR CHECK: verify the NeuroC reference and result]. However, these studies do not jointly address robustness verification and deployment on the target platform. Our work combines these perspectives by evaluating predefined 8-bit and 16-bit configurations, generating a fixed-point C implementation, deploying the selected model on an \AAOS{}-based Khadas VIM3, and investigating the cost of formally verifying its fixed-point behavior.

\subsection{Verification of Quantized Neural Networks}


Formal verification of quantized networks must model finite-width operations such as shifting and rounding. Mistry \emph{et al.} use MILP with bit-precise semantics, showing both the value and scalability challenges of solver-based QNN verification~\cite{Mistry2022MILP}; QEBVerif combines reachability analysis and MILP to bound quantization error~\cite{Zhang2023QEBVerif}; and CEG4N uses verification to guide quantization refinement and reports increased ESBMC timeouts for more complex networks and larger input regions~\cite{Matos2022CEG4N}. In contrast, our methodology fixes the quantization configuration and uses ESBMC to check certificate obligations linked to the generated fixed-point C implementation, targeting local classification robustness under a bounded input perturbation.

Overall, previous studies separately demonstrate efficient fixed-point CNN quantization, target-aware embedded inference, and formal analysis of quantized networks. Our study connects these concerns in a single automotive case study by comparing 8-bit and 16-bit quantization, verifying a selected fixed-point C implementation under a fixed perturbation region, and executing the same implementation on a VIM3 running \AAOS{}. The emphasis is therefore on characterizing the practical relationship between quantization accuracy, embedded deployment, and the time and memory requirements of bit-precise verification.

\subsection{Linear Bound Propagation}

Linear bound propagation is the main alternative to exact solving for real-valued networks. DeepPoly associates each neuron with one symbolic lower and one symbolic upper linear bound and back-substitutes these bounds through earlier layers~\cite{Singh2019DeepPoly}. CROWN computes similar bounds by propagating linear relaxations of the activation functions backward from the property to the input~\cite{3327345.3327402} [AUTHOR CHECK: confirm this is Zhang et al., CROWN, NeurIPS 2018]. Both approaches scale to networks far beyond the reach of exact procedures, but they reason about real-valued or floating-point arithmetic and do not model the rounding, clamping, and finite accumulators of a fixed-point implementation. Their bounds therefore cannot be used directly as guarantees for the generated C code. Our work reuses the CROWN back-substitution structure but performs it in exact integer arithmetic, includes the rounding of each rescaling step as an explicit error term, and treats the result only as a candidate certificate whose steps must be verified by ESBMC in harnesses linked to the deployed C implementation.


\subsection{Quantization Synthesis and the Implementation Gap}
Previous approaches integrate formal reasoning directly into the quantization process. CEG4N uses counterexamples to refine candidate quantization strategies and checks classification equivalence between the original and quantized networks~\cite{Matos2022CEG4N}. Quadapter computes backward layer preimages and uses MILP under-approximations to synthesize layer-wise quantization strategies that preserve properties such as robustness or backdoor-freeness~\cite{Zhang2024Quadapter}. These approaches demonstrate how verification can guide quantization selection. In contrast, our current methodology does not use formal verification to synthesize the quantization configuration. A single global 16-bit fixed-point configuration is fixed before verification, and 8-bit configurations are evaluated afterwards on the validation split for comparison. We then verify local robustness of the fixed 16-bit implementation under a fixed input-perturbation bound.



A separate challenge concerns the gap between the mathematical quantized model and its deployed implementation. Guarantees for an abstract quantized model may not transfer directly to its deployed implementation~\cite{Cordeiro2025NNPL}. We therefore check robustness certificates linked to an explicit fixed-point C implementation and execute the same C file natively on an \AAOS{}-based Khadas VIM3. This design relates the cost of implementation-level verification to the behavior of the deployed model. Table~\ref{tab:comparison} positions our study against representative approaches.

\begin{table}[h]
\centering
\caption{Positioning against representative approaches. C-level includes C abstractions. [AUTHOR CHECK: confirm the remaining feature classifications.]
}
\label{tab:comparison}
\scriptsize
\resizebox{\columnwidth}{!}{%
\begin{tabular}{lccccc}
\toprule
Approach & Synth. & Fixed-point & C-level & Contract & AAOS\\
 & quant. & arithmetic & verification & composition & deploy.\\
\midrule
DeepPoly~\cite{Singh2019DeepPoly} & -- & -- & -- & -- & --\\
QNNVerifier~\cite{Song2021QNNVerifier} & -- & \checkmark & \checkmark & -- & --\\
CEG4N~\cite{Matos2022CEG4N} & \checkmark & -- & \checkmark & -- & --\\
Quadapter~\cite{Zhang2024Quadapter} & \checkmark & -- & -- & \checkmark & --\\
CROWN~\cite{3327345.3327402} & -- & -- & -- & -- & --\\
\textbf{This paper} &  --  &\checkmark & \checkmark & \checkmark & \checkmark\\ 
\bottomrule
\end{tabular}
}%
\end{table}

\section{Deployment-Aware Methodology}
\label{sec:methodology}

\subsection{Problem Definition and Verification Boundary}
\label{sec:problem}

Let $f:\mathbb{R}^{m}\rightarrow\mathbb{R}^{K}$ be a trained floating-point classifier, let $x$ be a preprocessed image with label $c$, and let $\mathcal{X}$ be the valid input domain. We define the local region
\begin{equation}
\mathcal{B}_{\varepsilon}(x)=\{x'\in\mathcal{X}:\|x'-x\|_{\infty}\leq\varepsilon\}.
\label{eq:ball}
\end{equation}
The radius is expressed in raw byte units of the $32\times32$ RGB image, before normalization, and $\mathcal{X}$ contains only valid byte images. It models bounded pixel or sensor variation, not every physical image transformation.

The deployed network is an integer C program $g$. The input encoder $Q_{\mathrm{in}}$ maps the declared input to the integer representation consumed by $g$. The target property is
\begin{equation}
\forall x'\in\mathcal{B}_{\varepsilon}(x),\quad g_c(Q_{\mathrm{in}}(x'))>g_k(Q_{\mathrm{in}}(x'))\quad\forall k\neq c.
\label{eq:robustness}
\end{equation}
The strict inequality excludes ties. The verification boundary begins at the preprocessed tensor and includes input encoding, fixed-point layers, rescaling, activation, saturation, and the output comparison. Image acquisition, sign detection, camera drivers, and vehicle actuation are outside the boundary unless explicitly modeled.

\subsection{Fixed-Point Semantics}
\label{sec:synthesis}


A signed layer format is written $\langle Q,I,F\rangle$, where $Q=1+I+F$, $I$ excludes the sign bit, and $F$ is the fractional width. Integer $n$ represents $n\,2^{-F}$. Here Q7.8 denotes $\langle16,7,8\rangle$, with resolution $1/256$ and range $[-128,127.99609375]$, independently of the C storage type.

For the deployed configuration, all layers share Q7.8, and each affine layer accumulates into a signed \texttt{\_\_int128} value. The dot product is divided by 256 with round-half-away-from-zero, the quantized bias is added, and the result is clamped to the 16-bit range before any ReLU. Hidden layers use ReLU; the output layer does not. The certificate-linked C file stores parameters and internal activations as \texttt{int64\_t}; a compact variant uses \texttt{int16\_t}. The principal 8-bit comparison uses $\langle8,3,4\rangle$ (Q3.4) with the same rounding and saturation rules and \texttt{\_\_int128} accumulation. A Python-computed static bound excludes overflow of the deployed dot-product accumulators; this is distinct from the sampled ESBMC arithmetic-safety checks.

The board executes the reference C file identified by the certificate's SHA-256 and a separate compact variant. The harnesses encode certificate identities and arithmetic lemmas. A trusted Python binding step links them to the deployed parameter arrays, with sampled ESBMC corroboration. This conditional source correspondence does not prove compiler correctness, binary equivalence, or absence of hardware faults.


\subsection{Contracts and Composition}
For layer $l$, let $f_l^C$ denote the finite-width C computation and let $\mathcal{A}_l$ and $\mathcal{G}_l$ denote its admissible input and output regions. The layer obligations are
\begin{align}
Q_{\mathrm{in}}(\mathcal{B}_{\varepsilon}(x))&\subseteq\mathcal{A}_1, \\
\forall a\in\mathcal{A}_l:\; f_l^C(a)&\in\mathcal{G}_l, \\
\mathcal{G}_l&\subseteq\mathcal{A}_{l+1}\quad(l<L), \\
\mathcal{G}_L&\subseteq\{z:z_c>z_k,\;k\neq c\}.
\label{eq:contracts}
\end{align}
The third condition is essential: independently verified layer harnesses do not compose unless the output of one layer satisfies the assumptions of the next. Integer scale conversions are included in the checked computation. If hidden layers are decomposed into output-neuron blocks, every block is checked under the same layer format and input contract; all blocks must pass.

Interval bounds did not close the margins in the reported certified cases. We therefore use untrusted finders to propose per-neuron interval bounds and exact-integer CROWN bounds relating neurons or output margins to earlier layers~\cite{3327345.3327402}. ESBMC checks the resulting certificate in C, including integer relaxations, rounding terms, and the affine, ReLU, and rescaling steps. A fact may be used only after its associated harnesses are verified and it has one verified source. ESBMC analyzes the generated harnesses with fully unwound loops; additional arithmetic-safety checks cover sampled obligations. A region is \emph{certified} when the complete required certificate closes under the stated binding, composition, and arithmetic assumptions. Refutation requires a concrete violating network input; failed abstract margin checks alone do not establish one. Timeouts, memory exhaustion, unknown results, and unresolved required obligations are classified as \emph{inconclusive}.



\subsection{Native Deployment on \AAOS{}}
\label{sec:aaos}

The C backend is compiled into a statically linked AArch64 executable using GCC and run from the device shell. The target is a Khadas VIM3 with an Amlogic A311D system-on-chip. The board provides four \Aseven{} performance cores and two \Athree{} efficiency cores~\cite{KhadasVIM3}; AOSP documents the VIM3 as a reference board capable of running AOSP-based builds~\cite{AOSPReferenceBoards}. The experiment records the build fingerprint, kernel string, CPU information, compiler flags, and executable and data hashes.

The comparison uses one inference worker with affinity restricted to the \Aseven{} cluster (CPUs 2--5). The integer variants share fixed-point values and inputs, with different storage types. No warm-up inferences are discarded; paired \Athree{}/\Aseven{} and all-core comparisons were not measured.

\subsection{Evaluation Metrics}
The formal evaluation reports the number and fraction of certified, refuted, inconclusive, and misclassified image--radius pairs; CROWN finder and ESBMC time; peak ESBMC memory; and the number of ESBMC calls per layer. The denominator includes every declared attempt, including failures, timeouts, and resource aborts. As a scalability reference, we also report how far a monolithic ESBMC run on the C program progresses under a separately reported 600\,s timeout and \texttt{5g} memory limit.

The deployment evaluation reports classification accuracy of the float32 and fixed-point builds on the validation split, integer-output parity between the board and the host execution of the certified C file, median and 95th-percentile encoding-plus-inference latency, peak resident memory, and binary size. CPU utilization, frequency, temperature, and energy per inference were not measured.


\section{Experimental Configuration}
\label{sec:setup}

\subsection{Research Questions and Data Protocol}

The study evaluates the following questions:

% Research questions
\begin{tcolorbox}[
    colback=gray!10,
    colframe=gray!50,
    boxrule=0.5pt,
    arc=1mm,
    left=2mm,
    right=2mm,
    top=1mm,
    bottom=1mm
]
\begin{description}[leftmargin=0.55cm,labelsep=0.2em,itemsep=1pt,topsep=0pt]
  \item[RQ1:] Which GTSRB image--radius pairs satisfy the complete integer-C certification criterion?
  \item[RQ2:] What finder and ESBMC costs are required, and does decomposition bound individual query costs where monolithic BMC does not?
  \item[RQ3:] What latency, variability, and resource usage does the certified C implementation exhibit on \AAOS{}, compared with a float32 reference?
  \item[RQ4:] Does the deployed executable preserve the host integer outputs and the predictions of the quantized reference model?
\end{description}
\end{tcolorbox}

The benchmark is GTSRB, a single-image, 43-class traffic-sign classification dataset containing more than 50,000 images and realistic variation in illumination, weather, occlusion, scale, and rotation~\cite{Stallkamp2011GTSRB}. We use the official release: 39,209 training images and 12,630 test images. The training images are split by physical sign track (seed 2026, fraction 0.2) into 31,379 training and 7,830 validation images, so no track spans two partitions. Images are resized to $32\times32$ RGB by center-aligned nearest-neighbor sampling and normalized by $1/256$, the same map implemented by the exact byte encoder of the deployed C. Model selection used a predeclared 90\% validation-accuracy gate. Quantization-format search, finder development, and all board measurements use only the validation split. The selected model and its C implementation were evaluated once for test accuracy, after selection; the separate verification campaign used the draw described below.

The network has five layers: a $5\times5$ stride-2 convolution (6,144 outputs), two further convolutions (4,096 and 3,072 outputs), a learned $2\times2$ stride-2 depthwise convolution with a densely stored kernel (768 outputs), and a dense output layer (43 logits). It has 14,123 neurons and 307,211 stored parameter values. The float32 source reaches 91.32\% test accuracy. The deployed integer C reaches 91.35\% (balanced 87.33\%; worst class 30 at 44.7\%). Its logits are bit-exact with the Python integer reference on all 12,630 test images.

The radius is $\varepsilon = 1$ raw byte in $L_\infty$ over the 3,072 input bytes, i.e.\ $1/256$ after normalization. The property is local robustness: the target logit is strictly greater than every competing logit for every byte image in the ball. The verification sample is 80 test images drawn with seed 20260925. The draw, the deployment hash, all verifier source hashes, the resource limits, and the outcome taxonomy were pinned in a protocol file before the test verification campaign began. The formal denominator is all 80 drawn image--radius pairs, including images that the deployed C misclassifies.

\subsection{Hardware and Software}

Table~\ref{tab:setup} lists the configuration. Every benchmark binary is single-threaded and pinned with \texttt{taskset} to the \Aseven{} cluster (CPUs 2--5). Each build processes the 7,830 validation images once, checking outputs against the host, and then again in a 3-pass latency run (23,490 inferences). Resident memory comes from a separate 1-image, 1,000-pass run. No warm-up inferences are discarded. We report median and 95th-percentile latency; variability is reported as the p95$-$median spread and the maximum. The 99th percentile was not recorded. The sha256 of every binary and data file on the board was checked against the bundle manifest before the runs.

\begin{table}[h]
\centering
\caption{Experimental configuration used for deployment and formal verification.}
\label{tab:setup}
\scriptsize
\renewcommand{\arraystretch}{1.12}
\setlength{\tabcolsep}{4pt}
\begin{tabularx}{\columnwidth}{@{}>{\bfseries}p{0.26\columnwidth}X@{}}
\toprule
\textbf{Item} & \textbf{Configuration} \\
\midrule

Board &
Khadas VIM3, Amlogic A311D; 4$\times$ \Aseven{} and
2$\times$ \Athree{}; 3,744\,MiB reported RAM. \\

\addlinespace[1pt]
OS &
\AAOS{} \texttt{aosp\_yukawa\_car}, Android 16 (Baklava),
BP1A.250505.005.B1, \texttt{userdebug}; kernel
7.1.0-mainline with 4\,KiB pages. [AUTHOR CHECK: confirm automotive features, Android release/API level, and page size.] \\

\addlinespace[1pt]
Native toolchain &
GCC 13.3.0 \texttt{aarch64-linux-gnu} cross-compiler;
AArch64 ABI; statically linked; flags:
\texttt{-std=c11 -O2 -Wall -Wextra -Werror -static}. \\

\addlinespace[1pt]
Model &
5-layer GTSRB CNN with 14,123 neurons and 307,211 stored
parameters; global $\langle16,7,8\rangle$ format with signed 16-bit numerical range,
round-half-away, clamp before ReLU, and
\texttt{\_\_int128} accumulators. The certified C
implementation stores parameters as \texttt{int64\_t}. \\

\addlinespace[1pt]
Formal tools &
\ESBMC{} 7.11.0 (64-bit) with Z3 using profile
\texttt{paper-z3}; Bitwuzla also used for direct-C
baselines and auxiliary checks. Untrusted finders use interval propagation and
exact-integer CROWN back-substitution; no MILP is used in
this campaign. \\

\addlinespace[1pt]
Resource limits &
Box stage: 120\,s, \texttt{6g} per harness, 4 jobs, and 24-neuron blocks.
Chain check: 60\,s and \texttt{4g} per harness, with the job count
adjusted from $4\!\rightarrow\!12\!\rightarrow\!6\!\rightarrow\!4$
for resource management. Host free-memory guard: 6\,GiB. \\

\addlinespace[1pt]
Verification host &
x86-64, Linux 6.18 (WSL2), 20 GiB RAM. [AUTHOR CHECK: provide recorded host kernel/RAM metadata.] \\

\bottomrule
\end{tabularx}
\end{table}



\section{Results and Discussion}
\label{sec:results}

\subsection{Formal Certification (RQ1)}
A pair counts as certified only if two things hold:
\begin{itemize}
  \item every required ESBMC obligation of its certificate returns VERIFIED in the generated C harnesses, covering input non-vacuity, encoding, every layer, the rounding lemma, the ReLU and breakpoint checks, and the output margin;
  \item the dependency closure is complete: each fact has exactly one source, either an ESBMC-proved box block or an ESBMC-checked chain.
\end{itemize}

\begin{table*}[t]
\centering
\caption{Formal verification results for the 80 pre-registered GTSRB test pairs with $\varepsilon=1$ byte.}
\label{tab:formal}
\scriptsize
\renewcommand{\arraystretch}{1.15}
\setlength{\tabcolsep}{7pt}
\begin{tabular}{lrrrrrrr}
\toprule
& & \multicolumn{3}{c}{\textbf{Verification Outcome}}
& \multicolumn{2}{c}{\textbf{Inconclusive}}
& \\
\cmidrule(lr){3-5}
\cmidrule(lr){6-7}
\textbf{Run}
& \textbf{Pairs}
& \textbf{Miscl.}
& \textbf{Certified}
& \textbf{Refuted}
& \textbf{Finder}
& \textbf{Checker}
& \textbf{Certification Rate (95\% Wilson CI)} \\
\midrule

Primary (pre-registered)
& 80
& 8 (10.0\%)
& 1 (1.25\%)
& 0 (0.0\%)
& 54 (67.5\%)
& 17 (21.25\%)
& 1.25\% [0.2, 6.7] \\

With disclosed rerun
& 80
& 8 (10.0\%)
& 2 (2.50\%)
& 0 (0.0\%)
& 54 (67.5\%)
& 16 (20.0\%)
& 2.50\% [0.7, 8.7] \\

\bottomrule
\end{tabular}

\end{table*}

Table~\ref{tab:formal} reports one radius. ESBMC certified two test images:
\begin{itemize}
  \item image 06267 (class 40) with 48,871 CROWN-checker obligations;
  \item image 04083 (class 15) with 50,851 CROWN-checker obligations, in a disclosed rerun.
\end{itemize}
A validation development image was also certified, with 58,695/58,695 CROWN-checker obligations in the hardened run. These counts exclude box-stage calls; the image is outside the test sample.

The largest certified radius in the full-network test campaign is $\varepsilon = 1$ byte, the only radius attempted in that campaign.

No concrete network counterexample was established for the sample. Failed abstract box margins are closed by the completed CROWN certificates. After the rerun, the inconclusive pairs have two recorded causes:
\begin{itemize}
  \item \textbf{Relaxation gap (54 pairs).} The untrusted CROWN finder could not produce a positive margin lower bound for at least one competitor class (median minimum screen margin $-6{,}038.5$ over these 54 pairs). This is a property of the relaxation, not evidence of a reachable counterexample.
  \item \textbf{Resource abort (16 pairs).} The finder was positive, but the checker stopped because host free memory fell below its guard. The in-flight obligations were recorded as UNKNOWN or MEMOUT, never as FAILED.
\end{itemize}
If every finder-positive pair verified, the rate would be 18/80 = 22.5\% [14.7, 32.8]. This is the maximum attainable certification rate for this sample with the recorded finder outputs, not an achieved result.

The guarantee is local. If the complete proof discharges Eq.~\eqref{eq:contracts}, every input in the declared region is covered under the implemented integer semantics. No conclusion follows for inputs outside that region, and a non-certified pair is not shown to be non-robust.

\subsection{Verification Cost and Scalability (RQ2)}

\begin{table*}[t]
\centering
\caption{Verification cost: elapsed seconds, summed across CROWN calls; box-stage wall time is separate. CROWN counts exclude box calls; CROWN RSS is aggregate across jobs. Direct-C layer rows check interval containment.}
\label{tab:cost}
\scriptsize
\renewcommand{\arraystretch}{1.15}
\setlength{\tabcolsep}{7pt}

\begin{tabular}{lrrrrl}
\toprule
\textbf{Configuration}
& \textbf{Finder (s)}
& \textbf{ESBMC (s)}
& \textbf{Calls}
& \textbf{Peak RSS}
& \textbf{Outcome} \\
\midrule

\multicolumn{6}{l}{\textit{Direct-C verification (Q7.8; 600\,s, \texttt{5g})}} \\
\addlinespace[1pt]

Whole net, concrete input (Z3)
& -- & 600 & 1 & 3.0\,GB
& TIMEOUT in symex \\

L0: 24/6,144 outputs, 12 free bytes (Bitwuzla)
& -- & 186 & 1 & 4.7\,GB
& VERIFIED \\

L0: 24 outputs, 18 free bytes (Bitwuzla)
& -- & 270 & 1 & 4.9\,GB
& OUT\_OF\_MEMORY \\

L0: 48 outputs, concrete input (Bitwuzla)
& -- & 122 & 1 & 5.4\,GB
& OUT\_OF\_MEMORY \\

\midrule

\multicolumn{6}{l}{\textit{Decomposed verification}} \\
\addlinespace[1pt]

Box stage, including 589 blocks (test 06267)
& -- & 616 (wall) & 689 & --
& 605 VERIFIED; 84 abstract-margin FAILED \\

CROWN chains, test 06267
& 109 & 42,541 & 48,871 & 4,196\,MiB
& VERIFIED \\

CROWN chains, test 04083
& 34 + export & 73,714 & 50,851 & 8,635\,MiB
& VERIFIED \\

CROWN chains, validation 3652 (dev.)
& -- & 50,123 & 58,695 & 4,064\,MiB
& VERIFIED \\

\bottomrule
\end{tabular}


\end{table*}
[AUTHOR CHECK: recover test 04083's missing export time.]

The whole-network direct-C query times out on a concrete input. Q7.8 layer-0 interval checks verify 24 outputs with 12 free bytes; 18 free bytes, or 48 concrete-input outputs, exhaust memory. Decomposed certificates cover all 14,123 neurons over the declared ball. Successful CROWN stages use 11.8--20.5 hours of summed elapsed time and 49k--59k calls, excluding box stages and unsuccessful attempts. Chain and box harnesses have \texttt{4g} and \texttt{6g} limits, respectively.

Decomposition reduces individual query size, but it multiplies the number of calls. Above about 4k multiply terms, ESBMC cost is reported as 177--188\,$\mu$s per exported term [AUTHOR CHECK: identify the supporting sample and calculation]. This suggests proportionality within the sample. Layer 2 accounts for 56--59\% of summed elapsed time in the completed runs. Image 04083 partly ran at 12 jobs under different concurrency and memory conditions.

Decomposition also changes what is proved. The whole-network direct-C query states the classification property directly; the layer slices check interval containment. The decomposed pipeline proves per-block interval facts and per-chain affine and ReLU lemmas in generated C harnesses. The property then follows by composition, and that composition rests on two parts outside ESBMC:
\begin{itemize}
  \item the checked dependency closure;
  \item a binding step, which maps certificate layers to the deployed C arrays and which ESBMC corroborates only on sampled rows.
\end{itemize}
Both are part of the trust base.

Every timeout, memory-out, and unknown outcome remains in the denominator of Table~\ref{tab:formal}.

\subsection{CPU Execution on \AAOS{} (RQ3)}

The two QNN rows in Fig.~\ref{fig:vim3-deployment} use the same Q7.8 arithmetic.

\begin{figure}[h]
    \centering
    \includegraphics[width=\columnwidth]{BKP/figs/vim3_deployment_results_v2.pdf}
    
    \caption{VIM3 deployment: validation accuracy, encoding-plus-inference latency, and one-image process peak RSS for float32 and the two QNN storage variants. [AUTHOR CHECK: replace single-pass figure latencies with the reported three-pass values; confirm the figure path.]}
    \label{fig:vim3-deployment}
\end{figure}


\begin{itemize}
  \item The \emph{certified} row executes the reference C file linked to the checked certificates.
  \item The \emph{compact} row stores the same values as \texttt{int16\_t}. Its link to the certified file is measured (identical logits on 7,830 images), not proved.
\end{itemize}
The float32 row is a reference and is not formally verified.





\paragraph{Latency.}
The integer QNN is slower than float32 on the \Aseven{}: 1.44$\times$ for the certified build and 1.26$\times$ for the compact build, using the 3-pass medians. The cost of \texttt{\_\_int128} accumulation may contribute, but the generated code was not profiled. The 3-pass medians (31.67, 27.75, and 21.99\,ms) are within 0.20\,ms of the single-pass values, and p95 is within 0.2\,ms of the median. Maxima reach 59--69\,ms in the 3-pass runs; their cause is not identified because frequency and throttling were not recorded.

\paragraph{Memory.}
The compact build has the smallest one-image process peak RSS (1,260\,KiB vs 1,856\,KiB for float32) and executable-section total. The certified int64 build has the largest RSS (3,132\,KiB). RSS includes process overhead. Only the compact variant improves on float32's RSS; neither integer build improves latency.



\subsection{Numerical Fidelity (RQ4)}
On the \Aseven{}, both builds reproduce the host reference exactly on all 7,830 validation images:
\begin{itemize}
  \item the certified C matches all 43 int64 logits of every image;
  \item the compact build also matches them exactly;
  \item no encoder rejects occurred.
\end{itemize}
Host-to-\Aseven{} parity is therefore 100\% (7,830/7,830), with maximum absolute logit error 0.



These values characterize implementation fidelity and do not enlarge the certified input set. Applying a certificate to execution of the reference C source on the \Aseven{} assumes that the compiler, runtime, and hardware preserve the modeled semantics. The compact variant additionally relies on an unproved source-transformation equivalence.

\subsection{Answers to the Research Questions}
\textbf{RQ1.} The pre-registered primary result is one of 80 test pairs certified at $\varepsilon = 1$ byte (1.25\%, 95\% Wilson CI [0.2, 6.7]). The disclosed rerun raises this to two of 80 (2.5\%, [0.7, 8.7]). No pair was refuted. After the rerun, 54 pairs remain finder-inconclusive and 16 checker-inconclusive.

\textbf{RQ2.} In the Q7.8 direct-C pilot, 24 layer-0 outputs verify with 12 free bytes, while 48 outputs on a concrete input exhaust memory. Decomposition certifies local regions of the full 14,123-neuron network, with 11.8--20.5 hours of summed CROWN-check elapsed time and about 50k calls per image.

\textbf{RQ3.} On the \Aseven{}, the QNN is 1.26--1.44$\times$ slower than float32. Only the int16 compact build has lower measured process RSS than float32.

\textbf{RQ4.} Integer parity between host and \Aseven{} is exact on 7,830 of 7,830 images.

Average accuracy is not presented as evidence of robustness: the model is 91.35\% accurate on the test split, while 1.25\% of the sampled pairs are certified in the primary run and 2.5\% after the disclosed rerun.

\section{Limitations and Threats to Validity}
\label{sec:limitations}

\textbf{Internal validity.} Training used the training partition; model selection, format comparisons, and finder development used the validation partition; the selected model was evaluated once on the test split for accuracy and then used for the pre-registered verification draw. Deployment and verification artifacts, resource limits, and amendments were recorded, and primary and rerun results are reported separately.

\textbf{External and construct validity.} GTSRB is a single-image classification benchmark rather than a complete automotive perception stack; therefore, the results may not generalize to other tasks, platforms, deeper networks, or vehicle control. Local $L_\infty$ robustness at one byte models only bounded input variation and does not cover camera faults, distribution shift, or transformations outside the declared region. The 80-pair sample estimates this pipeline's certification rate, not the global robustness of the model.

\textbf{Formal scope.} ESBMC verifies the generated harnesses under stated assumptions and finite bounds. The trust base includes ESBMC and Z3, the certificate dependency closure, layer-to-C-array binding, the input-encoder contract, static accumulator-bound computation, and ESBMC interval-analysis settings. [AUTHOR CHECK: document completeness; the aggregator lacks independent obligation enumeration and record-count enforcement.] The proof does not establish compiler correctness, binary equivalence, hardware or operating-system correctness, global robustness, or vehicle-level safety. The compact int16 build is linked to the certified C implementation only by measured output equality.



\section{Conclusion}
\label{sec:conclusion}

This paper presented a deployment-aware study of fixed-point quantization, implementation-level verification, and embedded execution on \AAOS{}. The selected 16-bit Q7.8 model achieved 91.35\% test accuracy and reproduced the Python integer reference exactly on all test images, whereas the Q3.4 8-bit configuration reached only 80.75\% validation accuracy. At a one-byte $L_\infty$ radius, one of 80 pre-registered test image--radius pairs was certified in the primary run and two after a disclosed rerun; no pair was refuted. However, most remaining cases were inconclusive because of relaxation gaps or checker resource limits. Decomposition enabled local certification across the complete network under the stated trust assumptions, but the successful CROWN stages required roughly 49,000--59,000 ESBMC calls and 11.8--20.5 hours of summed per-call elapsed time. On the Cortex-A73, the QNN builds were slower than the float32 reference, although the compact int16 variant used less memory. These results demonstrate practical fixed-point deployment on \AAOS{} while identifying verification scalability as the principal barrier to broader implementation-level robustness guarantees.
